\section{Maxwell's Equations from Electromagnetic Tensor}\label{sec:appB}
Maxwell's electromagnetic tensor is defined as,
\begin{equation}
    F^{\mu\nu} = \partial^\mu A^\nu - \partial^\nu A^\mu
\end{equation}
where $A^\mu = (\Phi, \vb{A})$ is the four-potential. Using the four potential, we can write the electric and magnetic fields as
$$\vb{E} = -\nabla \Phi -\pdv{\vb{A}}{t}\qquad \vb{B} = \nabla\times \vb{A}$$
Let us now find some components of the electromagnetic tensor,
\begin{align*}
    F^{0i} =  \partial^0 A^i - \partial^i A^0 = \partial_0 A^i +\partial_i A^0 = -\qty[ - \dv{t} A^i - \partial_i A^0] = -E^i\\
\end{align*}
We also have the expression of the magnetic field as, \begin{align*}
    B^i = \varepsilon^{ijk}\partial^j A^k = \frac{1}{2}\qty[\varepsilon^{ijk}\partial^j A^k +\varepsilon^{ijk}\partial^j A^k] &= \frac{1}{2}\qty[\varepsilon^{ijk}\partial^j A^k -\varepsilon^{ikj}\partial^j A^k]\\  &=\frac{1}{2}\qty[\varepsilon^{ijk}\partial^j A^k -\varepsilon^{ijk}\partial^k A^j]\\ &= \frac{1}{2}\varepsilon^{ijk} \qty(\partial^j A^k-\partial^k A^j)\\
    &=\frac{1}{2}\varepsilon^{ijk}F^{jk}
\end{align*}
where in one step we introduced a minus by interchanging the Levi-Civita position and then in the next step, we renamed dummy indices $j \leftrightarrow k$. Thus, the explicit matrix representation of the electromagnetic tensor becomes:
\[
\qty[F^{\mu\nu}] =
\begin{pmatrix}
0 & -E^1 & -E^2 & -E^3 \\
E^1 & 0 & B^3 & -B^2 \\
E^2 & -B^3 & 0 & -B^1 \\
E^3 & B^2 & -B^1 & 0
\end{pmatrix}.
\]